\begin{titlepage}
\vspace*{22mm}
{\Huge Probability\\[4mm] and Statistics\par}
\vspace{8mm}
{\Large A complete course companion\par}
\vspace{10mm}
From uncertain events to evidence-based conclusions.
\vfill
English study notes organised into ten units for an introductory
probability and statistics course.

The explanations,
derivations and worked examples in these notes are original instructional
material, not a transcription of lectures or the listed textbooks.
\vspace{8mm}

\textbf{How to study:} read the intuition, reproduce the derivation,
cover the solution and redo the example, then answer the practice questions.
\end{titlepage}
\setcounter{tocdepth}{1}
\tableofcontents
\clearpage
\section*{Scope, prerequisites and notation}
\addcontentsline{toc}{section}{Scope, prerequisites and notation}
The ten units cover descriptive statistics, probability, distributions,
estimation, hypothesis testing and regression. Inference for variances,
chi-square tests of independence and $F$ distributions are included as well.

This is an introductory course companion; it is
not an encyclopedia of every branch of probability and statistics.
Prerequisites are elementary algebra, derivatives, definite integrals,
and the idea of a vector. When we differentiate a likelihood, we use
$d\log x/dx=1/x$. When we integrate a density, we are measuring area.

\textbf{Population versus sample.} A population is the collection or
data-generating process we want to understand. A sample is what we observe.
A parameter describes the population; a statistic is computed from a sample.
Probability starts with a model and predicts data. Statistics starts with
data and assesses what models or parameter values are plausible.

\begin{center}
\begin{tabular}{ll}
$P(A)$ & probability of event $A$\\
$X,\ x$ & random variable, observed value\\
$\mu,\sigma^2$ & population mean, population variance\\
$\bar X,\ S^2$ & random sample mean, sample variance\\
$\bar x,\ s^2$ & their observed numerical values\\
$\theta,\ \hat\theta$ & parameter, estimator or its estimate\\
$F_X,\ f_X,\ p_X$ & CDF, density, probability mass function\\
$q_p$ & lower-tail quantile: $P(W\leq q_p)=p$\\
$\alpha$ & significance level; $1-\alpha$ is confidence level
\end{tabular}
\end{center}
All quantiles in these notes are \emph{lower-tail} quantiles. Thus
$z_{0.975}\approx1.96$ and $t_{0.975,9}\approx2.262$.
Unless stated otherwise, a random sample means independent, identically
distributed observations (i.i.d.). All numerical examples are illustrative.

\clearpage
\section{Week 1: data, description and probability}
\subsection{Sampling and the questions data can answer}
A \textbf{simple random sample} gives every subset of the prescribed size
the same probability of selection. A \textbf{stratified sample} samples
within meaningful groups; it can improve precision and representation.
A \textbf{cluster sample} selects groups and observes units inside them;
observations in one cluster may be correlated. A systematic sample takes
every $k$th item after a random start; hidden periodicity can distort it.
A convenience sample is easy to collect but may not represent the population.

Random \emph{sampling} supports population generalization. Random
\emph{assignment} of treatments helps distinguish treatment effects from
confounding. An observational association alone does not establish causation.
Increasing sample size reduces random sampling error, but does not remove
selection bias, nonresponse bias, measurement error or confounding.

A categorical variable records labels. A quantitative variable is discrete
when its possible values are countable, and continuous when values range
over intervals. A measurement recorded to two decimals may still model
an underlying continuous quantity. Choose graphs accordingly: bar charts
for categories, histograms and boxplots for quantitative distributions,
scatterplots for two quantitative variables, and time plots for time order.
For unequal-width histogram bins, use density heights so that areas, not
heights alone, represent relative frequencies.

\subsection{Location and spread}
For observations $x_1,\ldots,x_n$,
\[
\bar x=\frac1n\sum_{i=1}^n x_i,\qquad
s^2=\frac{\sum_{i=1}^n(x_i-\bar x)^2}{n-1},\qquad s=\sqrt{s^2}.
\]
The mean balances all observations and is sensitive to outliers. The median
is the central ordered value (the average of the two central values for even
$n$), and is more resistant. The mode is a most frequent value. The range is
$\max x_i-\min x_i$; the interquartile range is $IQR=Q_3-Q_1$.
Quartile conventions differ between software packages, especially for small
samples; state the convention when it matters. A boxplot often flags values
outside $[Q_1-1.5IQR,Q_3+1.5IQR]$. A flag is a reason to investigate, not an
automatic reason to delete a measurement.

Variance is in squared units; standard deviation is in the original units.
The standardized value $z=(x-\bar x)/s$ describes distance from the sample
mean in sample standard deviations. If every observation changes to
$y_i=a+bx_i$, then $\bar y=a+b\bar x$ and $s_y=|b|s_x$.

\textbf{Worked example.} For $2,4,4,6,9$, the sum is 25 and $\bar x=5$.
The median is 4; squared deviations are $9,1,1,1,16$. Therefore
$s^2=28/4=7$ and $s\approx2.646$. The larger last value pulls the mean above
the median. Dividing by 5 would compute the variance of these five numbers
as a complete finite population; dividing by 4 estimates a population
variance from a sample. Unit 6 derives why $n-1$ appears.

\subsection{Events and probability axioms}
The sample space $\Omega$ contains possible outcomes. An event is a set of
outcomes. The complement $A^c$ means ``not $A$'', the union $A\cup B$ means
``$A$ or $B$ or both'', and the intersection $A\cap B$ means ``both''.
Disjoint events have empty intersection. The axioms are
\[
P(A)\geq0,\qquad P(\Omega)=1,\qquad
P\!\left(\bigcup_i A_i\right)=\sum_iP(A_i)
\]
for a countable sequence of pairwise disjoint events. They imply
\[
P(A^c)=1-P(A),\qquad
P(A\cup B)=P(A)+P(B)-P(A\cap B).
\]
The subtraction corrects double counting of the overlap. More generally,
$P(\bigcup_i A_i)\leq\sum_iP(A_i)$, the union bound.

Only when elementary outcomes are equally likely may we use
$P(A)=|A|/|\Omega|$. Counting formulas are tools for that model:
\[
n!=n(n-1)\cdots1,\quad
P(n,r)=\frac{n!}{(n-r)!},\quad
\binom nr=\frac{n!}{r!(n-r)!},\quad 0!=1.
\]
Use the product rule for successive choices. Use permutations when order
matters and combinations when it does not. With replacement, $r$ ordered
draws from $n$ possibilities have $n^r$ possible sequences.

\textbf{Worked example.} Choose 3 components without replacement from
10, of which 2 are defective. The probability of at least one defective is
\[
1-\frac{\binom83}{\binom{10}3}=1-\frac{56}{120}=\frac8{15}.
\]
The complement is simpler than splitting into exactly one and exactly two.

\textbf{Model-building habit.} State the population, how the sample was
obtained, what is random, and which outcomes are assumed equally likely.
A correct calculation cannot rescue an unjustified sampling model.

\clearpage
\section{Week 2: conditional probability and Bayes' rule}
\subsection{Conditioning changes the reference population}
For $P(B)>0$, define
\[
P(A\mid B)=\frac{P(A\cap B)}{P(B)}.
\]
Once $B$ is known, we restrict attention to $B$ and renormalize its
probability to one. Rearranging gives the multiplication rule
$P(A\cap B)=P(B)P(A\mid B)$. For a sequence of events,
\[
P(A\cap B\cap C)=P(A)P(B\mid A)P(C\mid A\cap B).
\]
In a probability tree, multiply along a path and add across disjoint paths.
Conditioning is not the same as intervention: knowing a patient received a
treatment is different from randomly assigning the treatment.

\subsection{Independence and common traps}
Events $A$ and $B$ are independent if $P(A\cap B)=P(A)P(B)$.
For $P(B)>0$, this is equivalent to $P(A\mid B)=P(A)$.
Disjoint events with positive probabilities are \emph{not} independent:
knowing one occurred rules out the other. Sampling without replacement
usually introduces dependence. Replacing an item and remixing can justify
independent draws when the experiment is otherwise unchanged.

Pairwise independence of three events is weaker than mutual independence.
For two fair independent coin tosses, let $A$ be ``first is heads'', $B$ be
``second is heads'', and $C$ be ``the tosses match''. Each pair has
intersection probability $1/4$, but
$P(A\cap B\cap C)=1/4\ne(1/2)^3$. Mutual independence requires factorization
for every subcollection, not just every pair.

\textbf{Worked example.} Two draws without replacement from a bag with
3 red and 2 blue balls give
\[
P(R_1\cap R_2)=\frac35\frac24=\frac3{10},\qquad
P(R_2\mid R_1)=\frac12.
\]
With replacement the first probability would be $(3/5)^2=9/25$.

\subsection{Total probability and Bayes' theorem}
Let $B_1,\ldots,B_k$ be a partition: disjoint, exhaustive events with
positive probabilities. Splitting event $A$ among the cases gives
\[
P(A)=\sum_{j=1}^k P(A\mid B_j)P(B_j).
\]
Substituting this into the conditional probability definition yields
\[
P(B_i\mid A)=
\frac{P(A\mid B_i)P(B_i)}{\sum_jP(A\mid B_j)P(B_j)}.
\]
The prior $P(B_i)$ represents uncertainty before observing $A$;
$P(A\mid B_i)$ is the likelihood of that observation under case $i$;
the posterior $P(B_i\mid A)$ incorporates the observation.

\textbf{Worked example: an automated alert.} Suppose 1\% of files are
malicious. A detector flags 95\% of malicious files and 5\% of benign files.
Then
\[
P(M\mid +)=\frac{0.95(0.01)}{0.95(0.01)+0.05(0.99)}
=\frac{0.0095}{0.059}\approx0.1610.
\]
Among 10,000 files, the model expects 95 true alerts and 495 false alerts.
Only $95/(95+495)$ of alerts are malicious. A high detection rate does not
by itself imply a high posterior probability; the base rate matters.

In odds form, Bayes' rule is
\[
\frac{P(M\mid+)}{P(M^c\mid+)}=
\frac{P(M)}{P(M^c)}\frac{P(+\mid M)}{P(+\mid M^c)}.
\]
This separates prior odds from the likelihood ratio supplied by evidence.
Do not reverse $P(A\mid B)$ and $P(B\mid A)$.

\clearpage
\section{Week 3: random variables and joint distributions}
\subsection{PMFs, densities and cumulative distributions}
A random variable maps each experimental outcome to a number. For a
discrete $X$, the probability mass function (PMF) satisfies
\[
p_X(x)=P(X=x)\geq0,\qquad \sum_xp_X(x)=1.
\]
For an absolutely continuous $X$, a probability density (PDF) satisfies
\[
f_X(x)\geq0,\quad \int_{-\infty}^{\infty}f_X(x)\,dx=1,\quad
P(a<X\leq b)=\int_a^b f_X(x)\,dx.
\]
A density can exceed 1: only its total area must be 1. At an individual
point, $P(X=x)=0$, so endpoints do not affect interval probabilities for a
continuous distribution. For a discrete variable, endpoints matter.

The cumulative distribution function works in both cases:
\[
F_X(x)=P(X\leq x),\qquad P(a<X\leq b)=F_X(b)-F_X(a).
\]
It is nondecreasing, right-continuous, and has limits 0 and 1 at the two
infinities. A PMF produces jumps; where a continuous CDF is differentiable,
$f_X=F_X'$. For a general distribution,
$P(X=x)=F_X(x)-F_X(x^-)$; mixed distributions may have both jumps and density.

\textbf{Worked example.} Let $f(x)=cx$ for $0<x<2$, zero otherwise.
Normalization gives $1=c\int_0^2x\,dx=2c$, hence $c=1/2$ and
\[
F(x)=\begin{cases}0,&x\leq0,\\x^2/4,&0<x<2,\\1,&x\geq2.\end{cases}
\]
Thus $P(1<X<1.5)=1.5^2/4-1/4=0.3125$. The median solves
$F(m)=1/2$, so $m=\sqrt2$, not the midpoint 1.

\subsection{Joint, marginal and conditional laws}
For two discrete variables, $p_{X,Y}(x,y)=P(X=x,Y=y)$. Marginalize by
summing over the variable you do not need:
\[
p_X(x)=\sum_y p_{X,Y}(x,y),\quad
p_{X\mid Y}(x\mid y)=\frac{p_{X,Y}(x,y)}{p_Y(y)}.
\]
For continuous variables, replace sums with integrals:
\[
f_X(x)=\int_{-\infty}^{\infty} f_{X,Y}(x,y)\,dy,\quad
f_{X\mid Y}(x\mid y)=\frac{f_{X,Y}(x,y)}{f_Y(y)}
\]
when $f_Y(y)>0$. This density definition handles conditioning even though
$P(Y=y)=0$. For an event region $R$,
$P((X,Y)\in R)=\iint_R f_{X,Y}(x,y)\,dx\,dy$.
Draw the support before choosing integration limits.

Random variables are independent when their joint law factorizes into
their marginal laws: $p_{X,Y}=p_Xp_Y$ or $f_{X,Y}=f_Xf_Y$ (almost everywhere).
Equivalently, the joint CDF is $F_{X,Y}(x,y)=F_X(x)F_Y(y)$ for all $x,y$.

\textbf{Worked example: triangular support.} Let $f_{X,Y}(x,y)=2$ on
$0<x<y<1$. Its total integral is
$\int_0^1\int_0^y2\,dx\,dy=1$. Therefore
\[
f_Y(y)=2y,\qquad f_X(x)=2(1-x),\qquad
f_{X\mid Y}(x\mid y)=\frac1y\quad(0<x<y).
\]
Given $Y=y$, $X$ is uniform between 0 and $y$. The variables are dependent:
the possible values of one are restricted by the other. Indeed,
$P(X<Y/2\mid Y=y)=1/2$, hence $P(X<Y/2)=1/2$.

\subsection{Transforming a random variable}
The safest general method is to find $F_Y(y)=P(g(X)\leq y)$ and differentiate.
If $g$ is one-to-one and differentiable with inverse $h$,
\[
f_Y(y)=f_X(h(y))|h'(y)|.
\]
The absolute derivative adjusts for stretching of the horizontal scale.
For a non-one-to-one transformation, sum contributions from all inverse
branches. If $X\sim U(0,1)$ and $Y=-\log X$, then for $y\geq0$,
$F_Y(y)=P(X\geq e^{-y})=1-e^{-y}$, giving density $e^{-y}$.
For independent continuous $X,Y$, the density of $Z=X+Y$ is the convolution
$f_Z(z)=\int f_X(x)f_Y(z-x)\,dx$ over the permitted support.

\clearpage
\section{Week 4: expectation, variance and discrete models}
\subsection{Expectation is a probability-weighted average}
Provided the relevant absolute moment is finite,
\[
E[g(X)]=\sum_xg(x)p_X(x)\quad\hbox{or}\quad
E[g(X)]=\int g(x)f_X(x)\,dx.
\]
Use this directly: you do not need the full distribution of $g(X)$ first.
For a pair, use the joint PMF or density. Linearity gives
$E[aX+bY+c]=aE[X]+bE[Y]+c$ even without independence.
In general, $E[g(X)]\ne g(E[X])$.
An indicator $I_A$ equals 1 on $A$ and 0 otherwise, so $E[I_A]=P(A)$.
Breaking counts into indicators is often the easiest way to find a mean.

Variance is the expected squared distance from the mean:
\[
\operatorname{Var}(X)=E[(X-\mu)^2]=E[X^2]-\mu^2,\qquad
\operatorname{Var}(aX+b)=a^2\operatorname{Var}(X).
\]
The identity follows by expanding $(X-\mu)^2$ and using $E[X]=\mu$.
For the density $f(x)=x/2$ on $(0,2)$ from Unit 3,
\[
E[X]=\tfrac12\int_0^2x^2\,dx=\tfrac43,\quad
E[X^2]=2,\quad \operatorname{Var}(X)=2-\tfrac{16}9=\tfrac29.
\]

\subsection{Covariance, correlation and sums}
\[
\operatorname{Cov}(X,Y)=E[(X-\mu_X)(Y-\mu_Y)]
=E[XY]-E[X]E[Y].
\]
When both variances are positive, correlation is the dimensionless quantity
$\rho=\operatorname{Cov}(X,Y)/(\sigma_X\sigma_Y)$, with $-1\leq\rho\leq1$.
It measures linear association, not every kind of dependence.
\[
\operatorname{Var}\!\left(\sum_i a_iX_i\right)
=\sum_i a_i^2\operatorname{Var}(X_i)
+2\sum_{i<j}a_ia_j\operatorname{Cov}(X_i,X_j).
\]
Independence implies zero covariance when moments exist. The reverse fails:
if $X\sim U(-1,1)$ and $Y=X^2$, symmetry gives $E[X]=E[X^3]=0$, so covariance
is zero, although $Y$ is determined by $X$.
For independent $X,Y$, $\operatorname{Var}(X-Y)=\operatorname{Var}(X)+
\operatorname{Var}(Y)$: the minus sign is squared.

Conditional expectation averages over a conditional distribution. Two
particularly useful identities are
\[
E[X]=E[E[X\mid Y]],
\]
\[
\operatorname{Var}(X)=E[\operatorname{Var}(X\mid Y)]
+\operatorname{Var}(E[X\mid Y]).
\]
The second splits total variation into variation within groups and
variation between group means. It follows by writing
$X-E[X]=(X-E[X\mid Y])+(E[X\mid Y]-E[X])$; the cross term has expectation zero.

\subsection{Markov and Chebyshev bounds}
For nonnegative $X$ and $a>0$, $E[X]\geq aP(X\geq a)$, hence
$P(X\geq a)\leq E[X]/a$ (Markov). Apply this to $(X-\mu)^2$:
\[
P(|X-\mu|\geq k\sigma)\leq\frac1{k^2}\quad(k>0,\ \sigma>0).
\]
This is Chebyshev's inequality. It requires finite variance but no bell
shape. At least $8/9$ of probability lies within $3\sigma$ in the corresponding
strict interval. For a normal model the coverage is about 99.73\%; do not
confuse a distribution-free lower bound with a normal-model calculation.

\subsection{Bernoulli and binomial}
For one success/failure trial, $X\sim\operatorname{Bernoulli}(p)$ has
$P(X=1)=p$, $P(X=0)=1-p$, $E[X]=p$, and variance $p(1-p)$ because $X^2=X$.
For $n$ independent trials with the same success probability,
\[
X\sim\operatorname{Bin}(n,p),\qquad
P(X=k)=\binom nk p^k(1-p)^{n-k},\quad k=0,\ldots,n.
\]
There are $\binom nk$ possible success positions, each with probability
$p^k(1-p)^{n-k}$. Writing $X=\sum_{i=1}^n I_i$ proves
$E[X]=np$ and $\operatorname{Var}(X)=np(1-p)$.

\textbf{Worked example.} In 10 independent tests with failure probability
0.1, the probability of exactly two failures is
$\binom{10}2(0.1)^2(0.9)^8\approx0.19371$. The probability of at least one
failure is $1-0.9^{10}\approx0.65132$. The expected number is 1, not a
guarantee that every batch contains exactly one.

\subsection{Poisson counts}
For event counts in a fixed exposure,
\[
X\sim\operatorname{Pois}(\lambda),\qquad
P(X=k)=e^{-\lambda}\frac{\lambda^k}{k!},\quad k=0,1,\ldots.
\]
Both mean and variance equal $\lambda$. For example,
$E[X]=e^{-\lambda}\lambda\sum_{j\geq0}\lambda^j/j!=\lambda$;
similarly $E[X(X-1)]=\lambda^2$, giving the variance.
Independent Poisson counts add: rates add too. If a process has rate $r$
per unit time, the count in length $t$ has mean $rt$ under the homogeneous
Poisson-process assumptions of constant rate and independent increments.

\textbf{Worked example.} At 3 requests per minute, the number in 20 seconds
has $\lambda=1$. Thus $P(X\geq2)=1-e^{-1}(1+1)\approx0.26424$.
The Poisson approximation to $\operatorname{Bin}(n,p)$ uses $\lambda=np$
when $n$ is large and $p$ is small; it is an approximation, not an identity.

\subsection{Related counting models}
These supplementary models clarify when binomial assumptions fail.
For sampling $n$ items without replacement from $N$ with $K$ successes,
\[
P(X=k)=\frac{\binom Kk\binom{N-K}{n-k}}{\binom Nn},\quad
E[X]=n\frac KN,
\]
\[
\operatorname{Var}(X)=n\frac KN\left(1-\frac KN\right)\frac{N-n}{N-1}.
\]
This is hypergeometric; only feasible $k$ values have nonzero probability.
The final factor is the finite-population correction. For the trial number
of the first success, a geometric variable on $1,2,\ldots$ has
$P(X=k)=(1-p)^{k-1}p$, mean $1/p$, variance $(1-p)/p^2$.
Some books count failures before the first success instead; state the support.

\clearpage
\section{Week 5: continuous models and the law of large numbers}
\subsection{Uniform distributions}
If all equal-length subintervals of $(a,b)$ have equal probability,
\[
X\sim U(a,b),\quad f(x)=\frac1{b-a},\quad
E[X]=\frac{a+b}2,\quad \operatorname{Var}(X)=\frac{(b-a)^2}{12}.
\]
The CDF inside the interval is $(x-a)/(b-a)$. Symmetry gives the mean;
integrating $(x-(a+b)/2)^2/(b-a)$ gives the variance.
If a bus arrives uniformly between 0 and 12 minutes from now,
$P(X>8)=4/12=1/3$ and the mean wait is 6 minutes.

\subsection{Exponential waiting times}
With \emph{rate} $\lambda>0$,
\[
f(x)=\lambda e^{-\lambda x},\quad F(x)=1-e^{-\lambda x},\quad x\geq0,
\]
\[
E[X]=1/\lambda,\qquad \operatorname{Var}(X)=1/\lambda^2.
\]
The CDF is zero below 0. Some sources use a scale $\theta=1/\lambda$;
check the parameterization before substituting numbers.
Integration by parts derives $E[X]=1/\lambda$ and $E[X^2]=2/\lambda^2$.
The survival function $S(t)=P(X>t)=e^{-\lambda t}$ gives
\[
P(X>s+t\mid X>s)=\frac{e^{-\lambda(s+t)}}{e^{-\lambda s}}
=e^{-\lambda t}.
\]
This memoryless property describes a constant hazard, not aging or wear-out.
In a homogeneous Poisson process, ``no arrival by $t$'' has probability
$e^{-rt}$, so the first waiting time is exponential with rate $r$.
At rate 2 per hour, the probability of waiting more than 45 minutes is
$e^{-2(3/4)}=e^{-1.5}\approx0.22313$.

\subsection{The normal distribution}
\[
X\sim N(\mu,\sigma^2),\quad
f(x)=\frac1{\sigma\sqrt{2\pi}}
\exp\!\left[-\frac{(x-\mu)^2}{2\sigma^2}\right],\quad \sigma>0.
\]
It is symmetric about $\mu$; the second parameter here is the variance.
Standardize using $Z=(X-\mu)/\sigma\sim N(0,1)$. If $\Phi$ is the standard
normal CDF,
\[
P(a<X<b)=\Phi((b-\mu)/\sigma)-\Phi((a-\mu)/\sigma).
\]
The CDF has no elementary antiderivative; use a table or numerical software.
Symmetry gives $\Phi(-z)=1-\Phi(z)$. About 68.27\%, 95.45\%, and 99.73\%
lie within one, two, and three standard deviations.

\textbf{Worked example.} If a dimension $X\sim N(100,15^2)$, then
$P(X>130)=1-\Phi(2)\approx0.02275$. Its 90th percentile is
$100+15z_{0.90}\approx119.22$.
An affine transform of a normal variable is normal. A linear combination
of independent normal variables is normal, with mean and variance obtained
from Unit 4. Normal marginals alone do not ensure that every combination is
normal unless appropriate joint normality holds.

For $X\sim\operatorname{Bin}(n,p)$, a normal approximation uses mean $np$
and variance $np(1-p)$ when both expected successes and failures are
sufficiently large. A useful screening guideline is both at least 10;
extreme tails may still be inaccurate. Apply a continuity correction:
\[
P(X\leq k)\approx
\Phi\!\left(\frac{k+0.5-np}{\sqrt{np(1-p)}}\right).
\]
The half-unit correction matches the continuous area to integer-valued bins.

\subsection{Law of large numbers (LLN)}
For i.i.d. observations with finite mean $\mu$, the weak LLN says
$\bar X_n\xrightarrow{P}\mu$: for every $\varepsilon>0$,
$P(|\bar X_n-\mu|>\varepsilon)\to0$.
Under the stronger finite-variance assumption, a short proof is
\[
E[\bar X_n]=\mu,\quad \operatorname{Var}(\bar X_n)=\sigma^2/n,
\quad P(|\bar X_n-\mu|\geq\varepsilon)
\leq\frac{\sigma^2}{n\varepsilon^2}\to0.
\]
The LLN explains why averages and relative frequencies stabilize. It does
not say individual observations approach the mean, and it does not say a
run of heads makes tails more likely next. Under independence, each next
coin toss retains its original probability. Unit 6's central limit theorem
adds an approximate \emph{shape and scale} for the remaining error.

\clearpage
\section{Week 6: sampling distributions and the CLT}
\subsection{A statistic has its own distribution}
Before data are collected, $\bar X$, $S^2$ and $\hat p$ are random variables.
Their sampling distributions describe repeated samples obtained using the
same design. A histogram of individual observations estimates a population
distribution; a histogram of averages from many repeated samples estimates
the sampling distribution of the mean. These are different objects.

If $X_1,\ldots,X_n$ are i.i.d. with mean $\mu$ and variance $\sigma^2$,
\[
E[\bar X]=\mu,\qquad \operatorname{Var}(\bar X)=\frac{\sigma^2}{n},
\qquad \operatorname{SD}(\bar X)=\frac{\sigma}{\sqrt n}.
\]
The standard deviation of an estimator is its \textbf{standard error} (SE).
The estimated SE of the mean is $s/\sqrt n$. Increasing $n$ by a factor of
four halves the SE; doubling $n$ does not halve it. For simple random
sampling without replacement from a finite population, the corresponding
SE includes a finite-population correction, with its exact form depending
on the convention used for population variance.

\subsection{Central limit theorem}
For i.i.d. observations with finite mean and finite positive variance,
\[
\frac{\bar X-\mu}{\sigma/\sqrt n}
\xrightarrow{d}N(0,1).
\]
Thus for large enough $n$, $\bar X\approx N(\mu,\sigma^2/n)$, and
$\sum_i X_i\approx N(n\mu,n\sigma^2)$. If the observations are themselves
normal, these distributions are exact for every $n$.
There is no universal ``$n\geq30$ makes everything normal'' rule: skewness,
heavy tails, rare events and dependence affect approximation quality.
An infinite-variance distribution need not obey this version of the CLT.
Dependence requires different conditions and often a different SE.

\textbf{Worked example.} A process has mean 50 and SD 12. For 36 independent
observations, the sample mean has SE 2. The CLT approximation is
\[
P(\bar X>54)\approx P(Z>2)=0.02275.
\]
For normal observations this is exact up to numerical rounding. Using 12
instead of 2 in the denominator would answer a different question about a
single observation.

For Bernoulli observations, $\hat p=X/n$ with $X\sim\operatorname{Bin}(n,p)$:
$E[\hat p]=p$ and $\operatorname{Var}(\hat p)=p(1-p)/n$. This connects
binomial probabilities with approximate proportion inference.

\subsection{Why sample variance divides by $n-1$}
The algebraic identity
\[
\sum_i(X_i-\bar X)^2=\sum_i(X_i-\mu)^2-n(\bar X-\mu)^2
\]
gives expected numerator $n\sigma^2-n(\sigma^2/n)=(n-1)\sigma^2$.
Therefore $S^2=\sum_i(X_i-\bar X)^2/(n-1)$ is unbiased for $\sigma^2$.
Estimating the mean uses one degree of freedom: the residuals sum to zero,
so only $n-1$ can vary freely. Unbiasedness of $S^2$ does not imply that
$S=\sqrt{S^2}$ is unbiased for $\sigma$.

\subsection{Chi-square, Student's t and F}
If $Z_1,\ldots,Z_\nu$ are independent standard normals, then
$V=\sum_iZ_i^2\sim\chi^2_\nu$. It has mean $\nu$, variance $2\nu$ and
nonnegative support. It is right-skewed for small $\nu$.
For a normal random sample,
\[
\frac{(n-1)S^2}{\sigma^2}\sim\chi^2_{n-1},\qquad
\bar X\ \text{and}\ S^2\ \text{are independent}.
\]
These exact results require normality. Unbiasedness of $S^2$ alone does not.

If $Z\sim N(0,1)$ and $V\sim\chi^2_\nu$ are independent,
\[
T=\frac{Z}{\sqrt{V/\nu}}\sim t_\nu.
\]
The $t$ distribution is symmetric with heavier tails than a normal. Its
tails account for uncertainty from estimating $\sigma$. As $\nu$ increases,
it approaches the standard normal. In a normal random sample,
$T=(\bar X-\mu)/(S/\sqrt n)\sim t_{n-1}$.

For independent $U\sim\chi^2_{\nu_1}$ and $V\sim\chi^2_{\nu_2}$,
\[
F=\frac{U/\nu_1}{V/\nu_2}\sim F_{\nu_1,\nu_2}.
\]
It is positive and generally right-skewed; numerator and denominator
degrees of freedom have different roles. Also $T_\nu^2\sim F_{1,\nu}$.
These distributions support variance inference and regression tests.

\clearpage
\section{Week 7: estimation and confidence intervals}
\subsection{What makes an estimator useful?}
An estimator $\hat\theta$ is a rule based on sample data. Its bias is
$E[\hat\theta]-\theta$. Its mean squared error decomposes as
\[
\operatorname{MSE}(\hat\theta)
=\operatorname{Var}(\hat\theta)+\operatorname{Bias}(\hat\theta)^2.
\]
Unbiasedness is useful but not sufficient: a slightly biased estimator may
have much smaller variance and MSE. Consistency means
$\hat\theta_n\xrightarrow{P}\theta$. Precision concerns sampling variability;
accuracy also depends on bias and model validity.

\subsection{Maximum likelihood estimation (MLE)}
For i.i.d. data with model $f(x;\theta)$,
\[
L(\theta;x_1,\ldots,x_n)=\prod_{i=1}^nf(x_i;\theta),\qquad
\ell(\theta)=\sum_i\log f(x_i;\theta).
\]
For discrete data replace density by PMF. Data are fixed when maximizing;
the parameter varies. A likelihood is not a probability distribution over
parameter values. Logarithms preserve the maximizer and turn products into
sums. Find stationary points, check curvature and admissible boundaries,
and consider whether support depends on the parameter.

\textbf{Bernoulli derivation.} Let $k=\sum_i x_i$. Then
\[
\ell(p)=k\log p+(n-k)\log(1-p),\qquad
\ell'(p)=\frac kp-\frac{n-k}{1-p}.
\]
Solving gives $\hat p=k/n$; the second derivative is negative at an interior
solution. If $k=0$ or $k=n$, the maximum is at the corresponding boundary.

\textbf{Exponential derivation.} For a rate $\lambda$,
$\ell(\lambda)=n\log\lambda-\lambda\sum_i x_i$; differentiating gives
$\hat\lambda=n/\sum_i x_i=1/\bar x$ when the sum is positive.
The MLE need not be unbiased: for $n>1$,
$E[\hat\lambda]=n\lambda/(n-1)$ in this model.

For normal data with both parameters unknown,
\[
\ell(\mu,\sigma^2)=-\frac n2\log(2\pi\sigma^2)
-\frac1{2\sigma^2}\sum_i(x_i-\mu)^2.
\]
Differentiation gives $\hat\mu=\bar x$ and
$\hat\sigma^2_{\rm MLE}=\sum_i(x_i-\bar x)^2/n$.
The MLE variance divides by $n$; the unbiased sample variance divides by
$n-1$. They answer different optimization or estimation criteria.

As a boundary example, if $X_i\sim U(0,\theta)$, then
$L(\theta)=\theta^{-n}1\{\theta\geq\max_i x_i\}$, so
$\hat\theta=\max_i x_i$. Blindly setting a derivative to zero would miss it.

\subsection{Confidence is a property of a procedure}
A $100(1-\alpha)\%$ confidence interval is a random interval constructed so
that, under the assumed sampling model, it covers the fixed parameter in
$100(1-\alpha)\%$ of repeated samples (exactly or approximately).
After observing the data, the resulting interval either contains the
parameter or does not. A frequentist 95\% interval does not assign a 95\%
posterior probability to the fixed parameter; that requires a Bayesian
model and prior. It also does not describe the middle 95\% of observations.

A general construction is estimate $\pm$ critical value $\times$ SE.
More confidence makes the interval wider. More data generally make it
narrower. These intervals do not correct sampling bias.

\subsection{Intervals for a mean}
If $\sigma$ is known and sampling is normal (or a CLT is adequate),
\[
\bar x\pm z_{1-\alpha/2}\frac{\sigma}{\sqrt n}.
\]
If $\sigma$ is unknown and data are normal,
\[
\bar x\pm t_{1-\alpha/2,n-1}\frac{s}{\sqrt n}.
\]
For large samples the latter is often a good approximation beyond exact
normality; small samples with strong skew or outliers need more care.
Both follow by inverting the central probability statement for the pivot.

\textbf{Worked example.} Suppose $n=16$, $\bar x=52$, $s=8$. A 95\% normal-model
interval for $\mu$ is
\[
52\pm t_{0.975,15}\frac8{4}=52\pm2.13145(2)
\approx[47.737,56.263].
\]
The half-width is 4.263, while the sample SD is 8; they measure different
things. With known/planning SD $\sigma$ and desired half-width $E$, choose
$n\geq(z_{1-\alpha/2}\sigma/E)^2$ and round up. A planning SD is an assumption,
not something known from the sample size formula itself.

\subsection{Intervals for a Bernoulli proportion}
The familiar Wald interval is
\[
\hat p\pm z_{1-\alpha/2}\sqrt{\frac{\hat p(1-\hat p)}n}.
\]
It can have poor coverage for small samples or proportions near 0 or 1,
and can extend beyond $[0,1]$. A more reliable general-purpose alternative
is the Wilson interval. Write $z=z_{1-\alpha/2}$; its center and half-width are
\[
c=\frac{\hat p+z^2/(2n)}{1+z^2/n},\qquad
h=\frac{z}{1+z^2/n}
\sqrt{\frac{\hat p(1-\hat p)}n+\frac{z^2}{4n^2}}.
\]
Report $[c-h,c+h]$. It comes from inverting the score test, whose standard
error uses the hypothesized $p$, rather than freezing it at $\hat p$.
For 60 successes in 100 trials, a 95\% Wilson interval is approximately
$[0.502,0.691]$; the Wald interval is approximately $[0.504,0.696]$.
For planning, $n\geq z^2p_*(1-p_*)/E^2$ with an anticipated $p_*$;
use $p_*=1/2$ when no credible prior estimate is available.

\subsection{Intervals for variance and a variance ratio}
For normal data, invert the chi-square pivot to obtain
\[
\left[\frac{(n-1)s^2}{\chi^2_{1-\alpha/2,n-1}},
\frac{(n-1)s^2}{\chi^2_{\alpha/2,n-1}}\right].
\]
The larger quantile goes in the lower endpoint's denominator. Square-root
the endpoints for a confidence interval for $\sigma$.
For $n=10$ and $s^2=4$, using quantiles $19.023$ and $2.700$ gives a 95\%
variance interval about $[1.892,13.333]$, which is asymmetric.

For independent normal samples, an interval for $\sigma_1^2/\sigma_2^2$ is
\[
\left[\frac{s_1^2/s_2^2}{F_{1-\alpha/2,n_1-1,n_2-1}},
\frac{s_1^2/s_2^2}{F_{\alpha/2,n_1-1,n_2-1}}\right].
\]
These variance procedures are sensitive to departures from normality.

\clearpage
\section{Week 8: hypothesis testing}
\subsection{The logic of a test}
Specify a null hypothesis $H_0$ and an alternative $H_1$ before inspecting
the result. A test statistic measures discrepancy from $H_0$ relative to
its sampling variation. The \textbf{p-value} is the probability, assuming
$H_0$ and its model assumptions, of a statistic at least as incompatible
with $H_0$ as the observed one, according to the stated alternative.
It is not $P(H_0\mid\text{data})$, and it is not the probability that the
result was ``due to chance''.

Choose a significance level $\alpha$, calculate the statistic and p-value,
and reject if $p\leq\alpha$. Otherwise, \emph{fail to reject}; this does not
prove equality. Report the estimated effect, interval, sample size and
assumptions as well as the decision. Statistical significance is not the
same as scientific or practical importance.

A Type I error rejects a true null. A level-$\alpha$ procedure controls its
probability at at most $\alpha$ under the null. A Type II error fails to reject
a false null; its probability $\beta$ depends on the actual alternative.
Power is $1-\beta$. Larger samples and larger true effects generally improve
power; lowering $\alpha$ at fixed sample size generally reduces power.

For a two-sided test of a scalar parameter, a matching equal-tail
$1-\alpha$ confidence interval excludes the null value exactly when that
test rejects (when constructed from the same procedure).
Repeated testing or selecting hypotheses after looking at the data changes
error rates. A Bonferroni adjustment uses $\alpha/m$ for each of $m$ tests
to control family-wise Type I error at at most $\alpha$.

\subsection{One normal population mean}
For $H_0:\mu=\mu_0$, known $\sigma$ gives
$Z=(\bar X-\mu_0)/(\sigma/\sqrt n)$, standard normal under $H_0$ for normal
data. For unknown $\sigma$, use
\[
T=\frac{\bar X-\mu_0}{S/\sqrt n}\sim t_{n-1}\quad\text{under }H_0.
\]
For $H_1:\mu\ne\mu_0$, $p=2[1-F_t(|t_{\rm obs}|)]$.
For $H_1:\mu>\mu_0$, $p=1-F_t(t_{\rm obs})$; for a lower alternative,
$p=F_t(t_{\rm obs})$. Do not halve a two-sided p-value after seeing the sign
unless the directional alternative was genuinely specified beforehand.

\textbf{Worked example.} With $n=25$, $\bar x=105$, $s=10$, test $\mu=100$
against $\mu\ne100$. Then $t=5/(10/5)=2.5$ with 24 degrees of freedom.
The two-sided p-value is about 0.0197; reject at 0.05. The matching 95\%
interval is $105\pm2.0639(2)=[100.872,109.128]$.
This supports a nonzero difference under the model; its practical relevance
depends on the units and application.

\subsection{Two independent means: Welch and pooled tests}
For independent groups and $H_0:\mu_1-\mu_2=\Delta_0$, the Welch statistic is
\[
t=\frac{\bar x_1-\bar x_2-\Delta_0}
{\sqrt{s_1^2/n_1+s_2^2/n_2}},
\]
with approximate degrees of freedom
\[
\nu=\frac{(s_1^2/n_1+s_2^2/n_2)^2}
{(s_1^2/n_1)^2/(n_1-1)+(s_2^2/n_2)^2/(n_2-1)}.
\]
A matching interval uses difference $\pm t_{1-\alpha/2,\nu}$ times that SE.
Welch does not assume equal population variances. Normality gives its usual
small-sample rationale; large samples permit broader approximations.

If equal variances are justified, the pooled estimate and SE are
\[
s_p^2=\frac{(n_1-1)s_1^2+(n_2-1)s_2^2}{n_1+n_2-2},\qquad
SE=s_p\sqrt{1/n_1+1/n_2},
\]
and the normal-model test has $n_1+n_2-2$ degrees of freedom.
Do not pool merely because a preliminary variance test was not significant.

\textbf{Worked example.} Let $n_1=20,\bar x_1=12,s_1=3$ and
$n_2=25,\bar x_2=10,s_2=4$. Under $\Delta_0=0$,
$SE=\sqrt{9/20+16/25}=\sqrt{1.09}\approx1.0440$,
$t\approx1.9157$ and $\nu\approx42.85$. The two-sided p-value is about
0.062. At 5\%, fail to reject. This does not establish that the means are equal.

\subsection{Paired observations}
For before/after data on the same subjects, or intentionally matched pairs,
define differences $D_i=X_i-Y_i$ and perform a one-sample test:
\[
t=\frac{\bar d-\Delta_0}{s_d/\sqrt n},\qquad df=n-1.
\]
Different pairs must be independent; within each pair dependence is expected.
Use the SD of the \emph{differences}, not the two group SDs as though groups
were independent. Pairing can substantially reduce variability if it is
scientifically meaningful. Choose a consistent subtraction order.

\subsection{Proportion and variance tests}
For $H_0:p=p_0$, a large-sample score statistic is
\[
z=\frac{\hat p-p_0}{\sqrt{p_0(1-p_0)/n}}.
\]
The null value belongs in its denominator. If expected counts are small,
use binomial probabilities for an exact procedure, specifying the convention
for a two-sided exact p-value. For two independent proportions under equality,
pool under the null: $\tilde p=(x_1+x_2)/(n_1+n_2)$ and
\[
z=\frac{\hat p_1-\hat p_2}
{\sqrt{\tilde p(1-\tilde p)(1/n_1+1/n_2)}}.
\]
An ordinary difference interval instead uses the unpooled estimated SE.

For normal observations and $H_0:\sigma^2=\sigma_0^2$, use
$Q=(n-1)S^2/\sigma_0^2\sim\chi^2_{n-1}$ under the null. A two-sided
equal-tail test rejects in either tail, with $\alpha/2$ in each.
For independent normal samples under equal variances,
$F=S_1^2/S_2^2\sim F_{n_1-1,n_2-1}$. Keep the stated ordering and use both
appropriate quantiles for a two-sided test.

\subsection{Chi-square tests for categorical counts}
For a goodness-of-fit test with observed counts $O_i$ and expected counts
$E_i=np_i$ under a specified categorical model,
\[
X^2=\sum_{i=1}^k\frac{(O_i-E_i)^2}{E_i}.
\]
Under suitable large-count conditions, compare to $\chi^2_{k-1}$ if all
category probabilities are specified in advance. If $m$ parameters are
estimated from these counts, the usual regular-model adjustment is
$df=k-1-m$. All expected counts must be positive. Small expected counts
can invalidate the approximation; use an appropriate exact or simulation
method, or scientifically justified category grouping.

For an $r\times c$ contingency table, independence predicts
\[
E_{ij}=\frac{(\text{row }i\text{ total})(\text{column }j\text{ total})}{n},
\quad X^2=\sum_{i,j}\frac{(O_{ij}-E_{ij})^2}{E_{ij}},
\quad df=(r-1)(c-1).
\]
Reject for a large statistic; the p-value is the upper chi-square tail.
Observations must be independent and each belongs to one cell. The
approximation is about expected counts, not observed counts alone.

\textbf{Worked example.} Two groups of 50 have pass/fail counts
$(30,20)$ and $(20,30)$. Under independence all four expected counts are 25.
Therefore $X^2=4(5^2/25)=4$, $df=1$, and the uncorrected Pearson p-value is
about 0.0455. At 5\% the data indicate an association. This calculation uses
no continuity correction; a corrected or exact test may give a different
p-value. Association does not by itself establish a causal group effect.

\subsection{A simple power calculation}
For a one-sided known-$\sigma$ test of $\mu=\mu_0$ versus $\mu>\mu_0$, reject
when $\bar X>\mu_0+z_{1-\alpha}\sigma/\sqrt n$.
At a true value $\mu_1$, normal-model power is
\[
1-\Phi\left(z_{1-\alpha}-\frac{\mu_1-\mu_0}{\sigma/\sqrt n}\right).
\]
For $\alpha=0.05$, $\sigma=10$, $n=25$ and $\mu_1-\mu_0=5$, power is
$1-\Phi(1.645-2.5)\approx0.804$. Planning power requires an effect size you
want to detect; there is no single power value for all alternatives.

\clearpage
\section{Week 9: simple linear regression}
\subsection{The model and what its coefficients mean}
For observations $(x_i,y_i)$, a simple linear regression model is
\[
Y_i=\beta_0+\beta_1x_i+\varepsilon_i,\qquad E[\varepsilon_i\mid x_i]=0.
\]
The conditional mean is $E[Y_i\mid x_i]=\beta_0+\beta_1x_i$. The slope is
the change in mean response per unit increase in $x$. The intercept is
the mean response at $x=0$; it may have little substantive meaning if zero
is outside the observed range. Errors describe departures from the unknown
population line; residuals $e_i=y_i-\hat y_i$ are computed after fitting.

For standard least-squares SE formulas assume uncorrelated errors with
constant variance $\sigma^2$. For exact small-sample $t$ and $F$ inference,
assume independent normal errors as well. Normality is not needed merely
to compute the line or for unbiasedness under the conditional mean model.

\subsection{Deriving the least-squares line}
Minimize $SSE(b_0,b_1)=\sum_i(y_i-b_0-b_1x_i)^2$.
Setting its two partial derivatives to zero produces the normal equations:
\[
\sum_i(y_i-b_0-b_1x_i)=0,\quad
\sum_i x_i(y_i-b_0-b_1x_i)=0.
\]
Define $S_{xx}=\sum_i(x_i-\bar x)^2$ and
$S_{xy}=\sum_i(x_i-\bar x)(y_i-\bar y)$. When $S_{xx}>0$,
\[
\hat\beta_1=\frac{S_{xy}}{S_{xx}},\qquad
\hat\beta_0=\bar y-\hat\beta_1\bar x,\qquad
\hat y_i=\hat\beta_0+\hat\beta_1x_i.
\]
The fitted line passes through $(\bar x,\bar y)$, and with an intercept
$\sum_i e_i=0$ and $\sum_i x_i e_i=0$. These are properties of the fit,
not evidence that the model is correct.

\subsection{Estimator properties and inference}
Because $\hat\beta_1=\sum_i(x_i-\bar x)Y_i/S_{xx}$ is a linear combination,
conditional mean zero errors give $E[\hat\beta_1]=\beta_1$. With uncorrelated,
equal-variance errors,
\[
\operatorname{Var}(\hat\beta_1)=\frac{\sigma^2}{S_{xx}},\qquad
\operatorname{Var}(\hat\beta_0)=\sigma^2\left(\frac1n+\frac{\bar x^2}{S_{xx}}\right),
\]
\[
\operatorname{Cov}(\hat\beta_0,\hat\beta_1)=-\frac{\sigma^2\bar x}{S_{xx}}.
\]
Under the Gauss--Markov assumptions, ordinary least squares is best
(minimum variance) among linear unbiased estimators. This claim does not
compare it to every possible biased or nonlinear estimator.

Estimate noise variance using $s^2=SSE/(n-2)$, because two coefficients have
been estimated. Then $SE(\hat\beta_1)=s/\sqrt{S_{xx}}$. Under the normal model,
\[
\frac{\hat\beta_1-\beta_{1,0}}{s/\sqrt{S_{xx}}}\sim t_{n-2}
\quad\text{under }H_0:\beta_1=\beta_{1,0}.
\]
A slope interval is $\hat\beta_1\pm t_{1-\alpha/2,n-2}s/\sqrt{S_{xx}}$.
An intercept interval uses its SE from the preceding variance formula.

\subsection{Explained variation and correlation}
Define $SST=\sum_i(y_i-\bar y)^2$ and
$SSR=\sum_i(\hat y_i-\bar y)^2$. With an intercept,
$SST=SSR+SSE$ and $R^2=1-SSE/SST$ when $SST>0$.
The value describes the fraction of observed response variation explained
by this fitted line. It does not prove causality or guarantee accurate
predictions on new data.

Pearson's sample correlation is $r=S_{xy}/\sqrt{S_{xx}S_{yy}}$.
In simple regression with an intercept, $R^2=r^2$ and
$\hat\beta_1=r(s_y/s_x)$. A normal-model test of zero slope is equivalent
to $F=(SSR/1)/(SSE/(n-2))$ with degrees of freedom $(1,n-2)$; $F=t^2$.
For i.i.d. bivariate normal data, a zero-correlation test uses
$t=r\sqrt{(n-2)/(1-r^2)}$ with $n-2$ degrees of freedom.

\subsection{Worked example from raw data}
Take $x=(1,2,3,4,5)$ and $y=(2,4,5,4,5)$. Then
$\bar x=3$, $\bar y=4$, $S_{xx}=10$, $S_{xy}=6$, and $S_{yy}=6$.
Hence $\hat\beta_1=0.6$, $\hat\beta_0=2.2$, and predictions are
$(2.8,3.4,4.0,4.6,5.2)$. Residuals are
$(-0.8,0.6,1.0,-0.6,-0.2)$, whose squares sum to 2.4.
Thus $s^2=2.4/3=0.8$, $R^2=1-2.4/6=0.6$, and
$r=\sqrt{0.6}\approx0.7746$.

The slope SE is $\sqrt{0.8/10}\approx0.28284$, giving
$t=0.6/0.28284\approx2.1213$ with only 3 degrees of freedom.
The 95\% slope interval is
$0.6\pm3.18245(0.28284)\approx[-0.300,1.500]$.
Despite $R^2=0.6$, this small sample gives substantial slope uncertainty.

\clearpage
\section{Week 10: diagnostics, prediction and validation}
\subsection{Residuals ask whether the model missed a pattern}
Plot residuals against fitted values and against each predictor.
A random-looking horizontal cloud is compatible with the mean/variance
assumptions; curvature suggests a missing nonlinear term, and a funnel
suggests nonconstant variance. Plot residuals against time or observation
order to look for dependence. A normal Q--Q plot compares ordered residuals
to normal quantiles; systematic departures or extreme tails challenge
normal-error inference. No diagnostic plot proves assumptions true.

An outlier has an unusual response relative to the model. A high-leverage
point has unusual predictor values. In simple regression its leverage is
\[
h_{ii}=\frac1n+\frac{(x_i-\bar x)^2}{S_{xx}},\qquad
\operatorname{Var}(e_i)=\sigma^2(1-h_{ii}).
\]
Thus residuals do not have identical variances even under a correct
equal-error-variance model. An internally standardized residual is
$r_i=e_i/[s\sqrt{1-h_{ii}}]$. A point is influential if removing it changes
the fitted model substantially. Cook's distance, for two fitted coefficients,
is
\[
D_i=\frac{e_i^2}{2s^2}\frac{h_{ii}}{(1-h_{ii})^2}.
\]
Large diagnostic values invite inspection of data quality and model
sensitivity. Do not remove points solely to achieve a desired p-value.

If curvature is substantive, consider a transformed predictor or response,
or a nonlinear mean function. If variance changes, weighted least squares
or heteroskedasticity-robust SEs may be appropriate. If observations are
dependent, use a model or variance estimator that respects the dependence.
These are extensions beyond the simple-regression core, not automatic fixes.

\subsection{Mean-response intervals versus prediction intervals}
At a predictor value $x_0$, the fitted mean is
$\hat y_0=\hat\beta_0+\hat\beta_1x_0$.
A normal-model confidence interval for the \emph{mean response} uses
\[
\hat y_0\pm t_{1-\alpha/2,n-2}s
\sqrt{\frac1n+\frac{(x_0-\bar x)^2}{S_{xx}}}.
\]
A prediction interval for one independent future response uses
\[
\hat y_0\pm t_{1-\alpha/2,n-2}s
\sqrt{1+\frac1n+\frac{(x_0-\bar x)^2}{S_{xx}}}.
\]
The extra 1 accounts for new observation noise, which remains even if the
mean line were known. Both intervals widen away from $\bar x$.
Extrapolation outside the sampled range additionally risks failure of the
model form; the displayed interval does not quantify that risk.

\textbf{Continue Unit 9's example.} At $x_0=3$, $\hat y_0=4$ and
$s=\sqrt{0.8}$. The mean-response SE is $\sqrt{0.8/5}=0.4$, so its
95\% interval is $4\pm3.18245(0.4)\approx[2.727,5.273]$.
The prediction SE is $\sqrt{0.8(1+1/5)}=\sqrt{0.96}\approx0.97980$,
giving $4\pm3.18245(0.97980)\approx[0.882,7.118]$.
The data are few; the wide intervals honestly reflect that.

\subsection{Validation measures performance on unseen data}
Training error is typically optimistic because fitting explicitly optimizes
performance on those observations. Separate fitting from evaluation.
Use a training set for fitting, a validation set or cross-validation for
choices, and a held-out test set for a final assessment when data permit.
For predictions $\hat y_i$ on held-out data,
\[
MAE=\frac1m\sum_i|y_i-\hat y_i|,\qquad
RMSE=\sqrt{\frac1m\sum_i(y_i-\hat y_i)^2}.
\]
RMSE penalizes large errors more strongly. Compare with a simple baseline,
such as a constant prediction fitted only on training data. Report units,
sample size and the data-splitting method. Test-set $R^2$ can be negative
if the model is worse than the reference mean used in its definition.

In $K$-fold cross-validation, split into $K$ folds, repeatedly fit on $K-1$
and evaluate on the remaining fold, then aggregate errors. All data-dependent
preprocessing, including scaling, feature selection and imputation, must
be learned within each training fold. Doing it once before splitting leaks
information. Keep grouped observations together; use forward-time splits
when predicting the future. Ordinary random folds can otherwise leak
subject or time information.

Model complexity trades bias against variance. An overly rigid model may
miss real structure; an overly flexible model may fit noise. Validation
helps choose useful complexity, but repeatedly tuning to the same test set
turns it into another validation set. Correlation is symmetric, while
regression predicts a particular response from specified predictors. Neither
alone turns observational data into causal evidence.

\clearpage
\section{A practical decision guide}
Start by identifying the scientific question, data type, sampling design,
target parameter and assumptions. Only then select a formula.
\begin{enumerate}
\item \textbf{Probability of an event:} draw its sample space; use counting
only for equally likely outcomes. For given evidence, condition; for a
mixture of cases, use total probability or Bayes.
\item \textbf{A count:} fixed independent trials suggest binomial; events
in an exposure suggest Poisson; sampling without replacement suggests
hypergeometric. State why the model fits the experiment.
\item \textbf{A measurement or wait:} consider normal, uniform or
exponential only when shape/support/process assumptions are appropriate.
\item \textbf{One mean:} known SD gives $z$; unknown SD gives $t$, with
exact normal-model inference or a suitable large-sample approximation.
\item \textbf{Two means:} matched observations require a paired analysis.
Independent groups generally use Welch; pooling needs equal-variance support.
\item \textbf{A proportion:} use a binomial model; Wilson intervals often
behave better than Wald. For tests, use the null-based SE or an exact method.
\item \textbf{A variance:} chi-square inference requires normal data;
variance ratios use $F$ for independent normal samples.
\item \textbf{Categorical association:} use a contingency table and assess
expected counts before a chi-square approximation.
\item \textbf{A numerical response versus a predictor:} fit regression,
inspect residuals and distinguish mean uncertainty from prediction noise.
\end{enumerate}

\textbf{Write a complete solution in five sentences.} (1) State the target
and model. (2) State the relevant assumptions. (3) Show the formula and
substituted values. (4) Report the numerical result with units and sensible
rounding. (5) Interpret it in the original context, including its limits.

\subsection*{Common errors to catch before submitting}
Do not multiply probabilities without a reason for independence. Do not
confuse mutually exclusive with independent. A PDF value is not a point
probability. Standard deviation is not standard error. A 95\% confidence
interval is not an interval containing 95\% of the observations. A p-value
is not the probability the null is true. ``Not significant'' does not mean
``no effect''. $R^2$ does not prove causation. Paired data are not two
independent samples. A large sample does not remove bias. Check supports,
units, tail directions, and whether the second normal parameter is variance
or standard deviation.

\clearpage
\section{Practice problems}
Try these before reading the solutions. All data below are invented for
practice; when a normal model is requested, treat its assumptions as given.
\begin{enumerate}
\item For $1,2,3,4,10$, find the mean, median and sample variance. Explain
which location measure better resists the last observation.
\item Draw 2 items without replacement from 8, of which 3 are defective.
Find the probability exactly one is defective.
\item Machines A and B produce 60\% and 40\% of items, with defect rates
1\% and 3\%. Given a defective item, find the probability it came from B.
\item A density is $f(x)=cx^2$ on $(0,1)$. Find $c$, the CDF there, the mean
and the variance.
\item A joint PMF on $(0,0),(0,1),(1,0),(1,1)$ is respectively
$0.2,0.3,0.1,0.4$. Find the two means, covariance and $P(X=1\mid Y=1)$.
Are the variables independent?
\item Independent $X,Y$ have means 2, 5 and variances 3, 4. Find the mean
and variance of $2X-Y+1$.
\item If $X\sim\operatorname{Bin}(8,0.25)$, find $E[X]$, its variance and
$P(X=0)$. If $Y\sim\operatorname{Pois}(2)$, find $P(Y\geq1)$.
\item An exponential wait has mean 4 minutes. Find its rate and
$P(X>6\mid X>2)$.
\item Measurements are i.i.d. $N(20,9)$. For $n=36$, find
$P(\bar X>21)$ exactly in terms of $\Phi$ and approximately numerically.
\item In 40 Bernoulli trials there are 14 successes. Derive the MLE of $p$.
For exponential observations $1,2,3$, find the MLE rate.
\item With normal data $n=9,\bar x=10,s=3$, find a 95\% interval for the mean.
Use $t_{0.975,8}=2.306$.
\item With normal data $n=16,\bar x=7,s=2$, test $\mu=6$ against $\mu\ne6$
at 5\%. Use $t_{0.975,15}=2.131$.
\item Five independent paired differences are $1,2,1,3,3$. Compute the
one-sample $t$ statistic for a zero mean difference and its degrees of freedom.
\item A $2\times2$ table has rows $(40,10)$ and $(30,20)$. Find expected
counts, the Pearson statistic and its degrees of freedom.
\item For the Unit 9 regression data, predict at $x=4$ and state whether
this is interpolation or extrapolation. Why is a prediction interval wider
than a confidence interval for the mean at that point?
\item A model scores very well on training data but poorly on a held-out
sample. Give two plausible explanations and one validation improvement.
\end{enumerate}

\clearpage
\section{Worked solutions to practice}
\begin{enumerate}
\item $\bar x=4$, median $=3$; squared deviations sum to
$9+4+1+0+36=50$, so $s^2=50/4=12.5$. The median is more resistant;
the last observation substantially raises the mean.
\item Choose one from each group: $\binom31\binom51/\binom82=15/28$.
The draws are dependent, which the combination calculation handles.
\item $P(D)=0.6(0.01)+0.4(0.03)=0.018$.
Bayes gives $P(B\mid D)=0.012/0.018=2/3$.
\item $c/3=1$, so $c=3$ and $F(x)=x^3$ on $(0,1)$.
$E[X]=3/4$, $E[X^2]=3/5$, hence variance $3/5-9/16=3/80$.
\item $E[X]=0.5$, $E[Y]=0.7$, $E[XY]=0.4$; covariance is
$0.4-0.5(0.7)=0.05$. Also $P(X=1\mid Y=1)=0.4/0.7=4/7$.
Not independent, since $P(X=1,Y=1)=0.4\ne0.35$.
\item Mean $=2(2)-5+1=0$. Independence gives variance
$2^2(3)+(-1)^2(4)=16$. The constant contributes no variance.
\item Binomial mean $=2$, variance $=8(0.25)(0.75)=1.5$, and
$P(X=0)=0.75^8\approx0.10011$. Poisson probability
$P(Y\geq1)=1-e^{-2}\approx0.86466$.
\item Rate $\lambda=1/4$ per minute. Conditional survival is
$e^{-(6-2)/4}=e^{-1}\approx0.36788$. Memorylessness concerns the remaining wait.
\item Population SD is 3, not 9. SE $=3/6=0.5$; the probability is
$1-\Phi((21-20)/0.5)=1-\Phi(2)\approx0.02275$.
\item $\ell(p)=14\log p+26\log(1-p)$; its derivative vanishes at
$p=14/40=0.35$, with negative second derivative. The exponential MLE is
$3/(1+2+3)=0.5$ per time unit.
\item SE $=3/3=1$, so $10\pm2.306=[7.694,12.306]$.
This is an interval for the population mean under the normal sampling model.
\item $t=(7-6)/(2/4)=2$, $df=15$. Since $|2|<2.131$, fail to reject.
The matching interval $7\pm2.131(0.5)$ contains 6.
\item $\bar d=2$ and $s_d^2=(1+0+1+1+1)/4=1$.
Thus $t=2/(1/\sqrt5)=2\sqrt5\approx4.4721$, $df=4$.
The use of differences preserves the paired design.
\item Column totals are 70 and 30, and row totals 50 each. Expected rows
are $(35,15)$ and $(35,15)$. Thus
$X^2=2(25/35+25/15)=100/21\approx4.7619$, $df=1$.
The upper-tail p-value is approximately 0.0291, without continuity correction.
\item $\hat y=2.2+0.6(4)=4.6$. It is interpolation within $[1,5]$.
The prediction interval includes future observation noise in addition to
uncertainty in the estimated conditional mean.
\item Possible causes include overfitting, distribution shift, and an
unrepresentative training sample. Use leakage-free cross-validation for
model choices and retain an untouched test set. If time or subjects create
dependence, split by time or subject rather than individual rows.
\end{enumerate}

\clearpage
\section{Course coverage and further study}
\begin{longtable}{p{.13\textwidth}p{.74\textwidth}}
\toprule
Sequence & Coverage in this companion\\\midrule
Week 1 & Unit 1: sampling, location, variability, data types,
graphs, sample spaces, events, equally likely outcomes.\\
Week 2 & Unit 2: conditional probability, independence, total probability,
Bayes' rule.\\
Week 3 & Unit 3: discrete and continuous random variables, joint,
marginal and conditional distributions.\\
Week 4 & Unit 4: expectation, variance, covariance, Chebyshev,
Bernoulli, binomial and Poisson.\\
Week 5 & Unit 5: uniform, exponential, normal, law of large numbers.\\
Week 6 & Unit 6: sample mean, CLT, sample variance, Student's $t$.\\
Week 7 & Unit 7: point estimation, MLE, mean and proportion intervals.\\
Week 8 & Unit 8: significance, errors, p-values, one- and two-mean tests.\\
Week 9 & Unit 9: least squares, fitted line, properties and coefficient inference.\\
Week 10 & Unit 10: residuals, checking, prediction, correlation, validation.\\
Extensions & Units 6--8: chi-square and $F$, variance inference,
chi-square independence testing.\\\bottomrule
\end{longtable}

\subsection*{Further reading}
The following textbooks offer additional explanations and exercises.
This companion does not reproduce their text or worked solutions.
\begin{enumerate}
\item Walpole, Myers, Myers, Ye. \emph{Probability and Statistics for Engineers
and Scientists}, 9th edition, Pearson.
\item Sheldon Ross. \emph{Introduction to Probability and Statistics for
Engineers and Scientists}, 5th edition, Elsevier.
\item Sheldon Ross. \emph{A First Course in Probability}, 10th edition,
Pearson.
\end{enumerate}

\textbf{Suggested ten-session revision cycle.} In each session, study the
matching unit, reproduce at least one derivation, solve an example without
looking, and explain the assumptions aloud. After Unit 5, revisit probability
and distribution selection. After Unit 10, mix the practice questions so
that choosing the method becomes part of the exercise. Revisit a mistake
until you can explain why the incorrect approach was tempting and why it fails.

\textbf{Study boundary.} This companion explains mathematics, not the
assessment requirements of a particular institution. Confirm course
requirements and examination logistics with your instructor.
